\section{Reflected checkers and bounded selection}

\subsection{Checker reflection}

A Boolean program is not useful in an object proof until its acceptance
condition is connected to a formula.

\begin{definition}[Reflected checker]
A reflected checker of arity \(k\) consists of a fixed PR symbol \(v\), a QF
formula \(A(\vec x)\), and object-iPRA derivations
\[
 v(\vec x)=1\to A(\vec x),
 \qquad
 A(\vec x)\to v(\vec x)=1,
 \qquad
 v(\vec x)=0\vee v(\vec x)=1.
\]
\end{definition}

The equivalence may be weakened if the compiler needs only soundness.  In
practice completeness is valuable because it lets mathematical existence
lemmas be converted into accepted-certificate existence.

\subsection{A checker algebra}

Fixed reflected checkers are closed under the finite logical operations needed
by certificate verification.

\begin{proposition}[Boolean closure]\label{prop:boolean-closure}
Suppose \(a,b\in\{0,1\}\) reflect formulas \(A,B\).  Then fixed PR terms can
be chosen to reflect
\(\neg A\), \(A\wedge B\), \(A\vee B\), and \(A\to B\).
\end{proposition}

\begin{proof}
Use
\[
 \neg a=1\mono a,
 \qquad
 a\wedge b=ab,
 \qquad
 a\vee b=1\mono((1\mono a)(1\mono b)).
\]
Implication is \(\neg a\vee b\).  Each truth table has four or fewer cases.
Those cases are object-iPRA disjunction eliminations from the supplied
Booleanity derivations, followed by semiring and monus equations already in the
arithmetic library.  Composing the existing soundness and completeness
implications gives reflection for the compound formula.
\end{proof}

Bounded quantifiers are finite folds.  For \(a(i,q)\in\{0,1\}\),
\[
  \forall i<N\ A(i,q)
  \quad\text{is reflected by}\quad
  \prod_{i<N}a(i,q)=1,
\]
and a bounded existential can be reflected by a nonzero bounded sum, a
bounded disjunction fold, or a least-search success test.  The repository's
bounded-sum and product libraries supply the required induction patterns.

\begin{theorem}[Bounded-formula compilation]\label{thm:bounded-compile}
Every formula built from reflected atomic equalities by intuitionistic
propositional connectives and numerically bounded quantifiers admits a fixed PR
Boolean checker with object-iPRA soundness and completeness.
\end{theorem}

\begin{proof}
Induct on the bounded formula.  Atomic equalities use \(\mathsf{eqb}\).
Propositional cases use \cref{prop:boolean-closure}.  For a bounded universal,
use a primitive-recursive product and prove by QF induction on the bound that
the product is one exactly when every previous factor is one.  For a bounded
existential, either dualize the universal or use the least-search specification.
All quantification introduced by the intended reading is removed from the
induction matrix: the running fold equation is QF.  The final bounded formula
is recovered by ordinary quantifier rules after the induction.
\end{proof}

\subsection{Selection from an explicit witness}

We now spell out the arithmetic reduction used in
\cref{thm:fqnf-compile}.  Let \(V(c,q)\) be Boolean and define
\(F(c,q)=1\mono V(c,q)\).  The checked search term
\(\mu_F(B,q)\) is the least \(c<B\) satisfying \(F(c,q)=0\), or \(B\) if no
such code exists.

\begin{theorem}[Least accepted certificate]\label{thm:least-cert}
From object proofs
\[
 w<B(q),\qquad V(w,q)=1,
\]
iPRA derives
\[
 \mu_F(B(q),q)<B(q),
 \qquad V(\mu_F(B(q),q),q)=1,
 \qquad \mu_F(B(q),q)\le w.
\]
\end{theorem}

\begin{proof}
The equation \(V(w,q)=1\) gives \(F(w,q)=0\).  Apply the repository theorem
\pathref{Search.search_found}; it returns a zero below the bound, no later
than \(w\).  Its zero equation is \(F(\mu,q)=0\).  By Booleanity, if
\(V(\mu,q)=0\) then \(F(\mu,q)=1\), contradicting the zero equation via
successor separation.  Therefore \(V(\mu,q)=1\).
\end{proof}

\subsection{Selection from existential production}

\begin{theorem}[Witness erasure]\label{thm:witness-erasure}
If iPRA proves
\[
  \exists w\,[w<B(q)\wedge V(w,q)=1],
\]
then it proves the first two conclusions of \cref{thm:least-cert} for the fixed
term \(\mu_F(B(q),q)\).
\end{theorem}

\begin{proof}
Use intuitionistic existential elimination.  In the temporary witness context,
apply \cref{thm:least-cert}.  The selected term depends on \(B,V,q\), not on the
temporary witness, so the conclusion is invariant under discharge.
\end{proof}

\subsection{Canonicality is optional}

The least-code convention yields a deterministic representative, but the
mathematics need not depend on it.  If every two accepted codes produce the
same or sufficiently close observable, later proofs should use that stability
theorem rather than properties peculiar to the numerical least code.  This
keeps the interface robust under changes of encoding and avoids introducing
spurious order into quotient constructions.
